% ============================================================
%  Erasure Is Deformation
%  Repair-Repertoire Reduction and the Geometry of Memory Loss
% ============================================================
%  Engine: LuaLaTeX
%  Font:   TeX Gyre Pagella
%  Sequel to: Observability Is Restorability (Flyxion, 2026)
% ============================================================

\documentclass[12pt, oneside]{article}

% ── Fonts & encoding ────────────────────────────────────────
\usepackage{fontspec}
\setmainfont{TeX Gyre Pagella}
\usepackage{unicode-math}
\setmathfont{TeX Gyre Pagella Math}

% ── Layout ──────────────────────────────────────────────────
\usepackage[
  top=1.25in, bottom=1.25in,
  left=1.35in, right=1.35in
]{geometry}
\usepackage{setspace}
\setstretch{1.15}
\setlength{\parskip}{0.55em}
\setlength{\parindent}{1.4em}

% ── Mathematics ─────────────────────────────────────────────
\usepackage{amsmath, amsthm, amssymb}
\usepackage{mathtools}

% Theorem environments
\theoremstyle{plain}
\newtheorem{theorem}{Theorem}[section]
\newtheorem{proposition}[theorem]{Proposition}
\newtheorem{lemma}[theorem]{Lemma}
\newtheorem{corollary}[theorem]{Corollary}

\theoremstyle{definition}
\newtheorem{definition}[theorem]{Definition}
\newtheorem{example}[theorem]{Example}
\newtheorem{axiom}[theorem]{Axiom}

\theoremstyle{remark}
\newtheorem{remark}[theorem]{Remark}

% ── Notation shorthands ─────────────────────────────────────
\DeclareMathOperator{\Rest}{Rest}
\DeclareMathOperator{\Reach}{Reach}
\DeclareMathOperator{\Obs}{Obs}
\newcommand{\CR}{C_{\!R}}
\newcommand{\AO}{\mathcal{A}}
\newcommand{\At}{\mathcal{A}_t}
\newcommand{\Azero}{\mathcal{A}_0}
\newcommand{\Rt}{\mathcal{R}_t}
\newcommand{\Rzero}{\mathcal{R}_0}
\newcommand{\Restzt}{\Rest(\Rt)}
\newcommand{\Restzero}{\Rest(\Rzero)}
\newcommand{\iotat}{\iota_t}
\newcommand{\Et}{E_t}
\newcommand{\gt}{g_t}
\newcommand{\gzero}{g_0}
\newcommand{\dRt}{d_R^{(t)}}
\newcommand{\IntMat}{\mathbf{I}}
\newcommand{\Hyster}{H}
\newcommand{\Mem}{\mathcal{M}}
\newcommand{\Dsub}{\mathcal{D}}
\newcommand{\Xsub}{X_S}
\newcommand{\bA}{\partial\mathcal{A}}

% ── Section & title styling ─────────────────────────────────
\usepackage{titlesec}
\titleformat{\section}{\normalfont\large\bfseries}{{\thesection}.}{0.6em}{}
\titleformat{\subsection}{\normalfont\normalsize\bfseries}%
  {{\thesubsection}.}{0.5em}{}

% ── Hyperref (load last) ────────────────────────────────────
\usepackage[
  colorlinks=true,
  linkcolor=black,
  citecolor=black,
  urlcolor=black
]{hyperref}

% ============================================================
\begin{document}

% ── Title block ─────────────────────────────────────────────
\begin{center}
  {\LARGE\bfseries Erasure Is Deformation}\\[0.9em]
  {\large Repair-Repertoire Reduction and the Geometry of Memory Loss}\\[1.5em]
  {\normalsize Flyxion}\\[0.3em]
  {\small\itshape Independent Researcher}\\[1.6em]
  {\small\today}
\end{center}

\vspace{0.5em}

% ── Abstract ────────────────────────────────────────────────
\begin{abstract}
\noindent
Memory is conventionally modelled as stored information.
This paper proposes a different account: memory is repair capacity,
and forgetting is the geometric deformation that results when repair
operators are removed.
We formalise erasure as reduction of a repair repertoire
$\Rzero \to \Rt \subseteq \Rzero$ and study its consequences for
the observer-relative admissibility manifold $\AO(\mathcal{R})$.
The Admissibility Contraction Theorem establishes that
$\At \subseteq \Azero$ with strict inclusion whenever distinctions
are lost.
The Erasure Deformation Theorem introduces the pullback deformation
tensor
$\Et = \iotat^* \gzero - \gt$,
where $\iotat : \Restzt \hookrightarrow \Restzero$ is the inclusion
of surviving repair domains, and proves that $|\Et| \to \infty$
along sequences approaching the irrestorability boundary.
Erasure therefore generates geometric singularities, not merely
information gaps.
The hysteresis functional $\Hyster(d) = \CR^{\mathrm{restore}}(d)
- \CR^{\mathrm{destroy}}(d)$ characterises irreversibility:
$\Hyster(d) = \infty$ is equivalent to permanent forgetting.
Topological memory is defined via the homology groups
$H_*(\At)$; component fragmentation, hole creation, and homology
collapse correspond to distinct modes of memory loss.
Applications are developed for neural injury, archive destruction,
language death, institutional collapse, and scientific knowledge loss.
The paper is a sequel to \emph{Observability Is Restorability}
\cite{Flyxion2026a}, extending that framework from the question of
what is observable to the question of what happens when observability
itself degrades.
\end{abstract}

\vspace{1em}
\hrule
\vspace{1.6em}

% ============================================================
\section{Memory, Erasure, and Deformation}
\label{sec:intro}
% ============================================================

The standard account of memory is a storage account.
Memories are records inscribed in a medium, and forgetting is the
degradation or loss of those records.
On this picture, memory loss is subtraction: what was present is
removed.

This paper begins from a different premise, developed formally in
the companion paper \emph{Observability Is Restorability}
\cite{Flyxion2026a}: memory is not storage but repair capacity.
To remember a distinction is to possess a repair path by which that
distinction can be restored after perturbation.
Forgetting is not the removal of a stored item but the loss of the
repair operators by which the item could be reconstructed.

If this is correct, then forgetting has geometric consequences that
the storage account cannot capture.
When repair operators are removed, the set of restorable distinctions
contracts, the observer-relative admissibility manifold deforms, and
the metric structure on surviving repair domains is distorted.
Erasure is not subtraction.
It is deformation.

The paper makes this claim precise.
We define erasure as reduction of a repair repertoire, study the
geometry of the resulting deformation, characterise irreversibility
via a hysteresis functional, and analyse the topological consequences
of progressive operator loss.
The framework applies equally to neural memory, institutional
knowledge, scientific technique, linguistic structure, and any other
system whose capacity for distinction-restoration can be degraded.

The argument proceeds as follows.
Section~\ref{sec:repair} reviews repair repertoires, restorable
distinctions, and the admissibility manifold from
\cite{Flyxion2026a}.
Section~\ref{sec:erasure} defines erasure events, the erasure
semigroup, and the lattice of sub-repertoires.
Section~\ref{sec:contraction} proves the Admissibility Contraction
Theorem.
Section~\ref{sec:deformation} introduces the pullback deformation
tensor and proves the Erasure Deformation Theorem.
Section~\ref{sec:hysteresis} develops the hysteresis functional and
the permanent forgetting condition.
Section~\ref{sec:topology} defines topological memory via homology
groups of the admissibility manifold.
Sections~\ref{sec:bio}--\ref{sec:science} develop applications.
Section~\ref{sec:consequences} draws consequences and identifies open
problems.


% ============================================================
\section{Repair Repertoires and Admissibility}
\label{sec:repair}
% ============================================================

We recall the essential definitions from \cite{Flyxion2026a}.
A \emph{substrate} $S = (\Xsub, \mathcal{P}, \Dsub)$ consists of a
locally compact Hausdorff state space $\Xsub$, a perturbation
semigroup $\mathcal{P}$, and a closed distinction relation
$\Dsub \subseteq \Xsub \times \Xsub$.
A \emph{repair repertoire} $\mathcal{R} \subseteq C(\Xsub, \Xsub)$
is a sequentially compact set of continuous maps equipped with the
compact-open topology.
The \emph{minimal repair cost} of distinction $d$ under repertoire
$\mathcal{R}$ is
\[
  \CR(d, \mathcal{R})
  \;=\;
  \sup_{p \in \mathcal{P}_d}
  \inf_{\gamma \in \Gamma_d^p(\mathcal{R})}
  \mathrm{cost}(\gamma),
\]
and the set of \emph{finitely restorable distinctions} is
\[
  \Rest(\mathcal{R})
  \;=\;
  \bigl\{d \in \Dsub \;\big|\; \CR(d, \mathcal{R}) < +\infty\bigr\}.
\]
The \emph{repair pseudometric} on $\Rest(\mathcal{R})$ is
\[
  d_R(d_i, d_j)
  \;=\;
  \inf_{\gamma : d_i \leadsto d_j}
  \int_\gamma c\, \mathrm{d}s,
\]
which defines a metric $g_{\mathcal{R}}$ on $\Rest(\mathcal{R})$
in the reversible case.
The \emph{observer-relative admissibility manifold} is
\[
  \AO(\mathcal{R})
  \;=\;
  \overline{\bigl\{
    x \in \Xsub \;\big|\;
    \forall\, d \in \Rest(\mathcal{R}),\;
    \CR(d, \mathcal{R}) < +\infty \text{ from } x
  \bigr\}}.
\]
The Repair Boundary Theorem \cite{Flyxion2026a} states that
$x \in \partial\AO(\mathcal{R})$ if and only if
$\CR(d, \mathcal{R}) \to +\infty$ for some $d \in \Rest(\mathcal{R})$
as the system approaches $x$.

\begin{remark}
  Throughout this paper, \emph{memory} denotes $\Rest(\mathcal{R})$
  together with its metric structure $g_{\mathcal{R}}$.
  This is not a storage register but a geometric object: the
  metrised space of distinctions the observer can maintain.
  Forgetting is change of this object.
\end{remark}

\begin{remark}[Observer persistence]
  Throughout this paper we assume that the observer persists across
  the erasure path: $\Rzero \to \Rt$ is a reduction belonging to
  the same observer, not a transition to a different observer.
  This assumption is non-trivial.
  Severe forgetting --- traumatic brain injury, late-stage
  neurodegeneration, total institutional collapse --- may change not
  only what the observer can restore but which distinctions the
  observer constitutes.
  In those cases $(\AO(\Rzero), \gzero)$ and $(\At, \gt)$ may be
  geometries of different observers rather than different states of
  the same one.
  The question of whether observer identity is itself defined by a
  path in the category of admissibility manifolds --- so that
  forgetting is \emph{observer deformation} rather than memory
  deformation --- is left for future work.
  The current framework is exactly correct for the regime in which
  identity persists: partial or moderate erasure where the surviving
  repair domain is non-empty and the observer remains recognisably
  continuous.
\end{remark}

\begin{remark}[The invariant making comparison meaningful]
  The reviewer of an earlier draft raised the question: if both the
  observer and the distinction set change under erasure, what
  invariant makes the comparison $\Et = \iotat^* \gzero - \gt$
  meaningful?
  The answer is the inclusion $\iotat : \Restzt \hookrightarrow
  \Restzero$.
  The pullback construction exists precisely because surviving
  distinctions are identified with their original representatives
  via $\iotat$.
  The deformation tensor measures how the geometry of those surviving
  distinctions has changed, relative to the geometry they had before
  erasure.
  No universal distinction space $\mathcal{D}_\infty$ is needed and
  none is assumed; the entire framework is observer-relative, which
  is the correct register for a repair-theoretic account.
  Introducing a universal distinction space would move the theory
  toward latent-space realism, which the parent paper
  \cite{Flyxion2026a} was designed to avoid.
\end{remark}


% ============================================================
\section{Erasure: Operators, Events, and the Sub-Repertoire Lattice}
\label{sec:erasure}
% ============================================================

\begin{definition}[Erasure event]
  An \emph{erasure event} at time $t$ is a map
  $\varepsilon_t : \mathcal{R} \mapsto \mathcal{R}'$ with
  $\mathcal{R}' \subseteq \mathcal{R}$, removing one or more repair
  operators from the available repertoire.
  An erasure event is \emph{local} if it removes operators acting
  on a bounded region of $\Xsub$; \emph{global} if it removes
  operators whose domain is all of $\Xsub$.
\end{definition}

\begin{definition}[Erasure path]
  An \emph{erasure path} is a decreasing family of repertoires
  $(\Rt)_{t \in [0,T]}$ with $\mathcal{R}_s \subseteq \mathcal{R}_t$
  for all $s \geq t$, and $\mathcal{R}_0$ the initial repertoire.
  The family is parameterised by $t$ increasing in the direction of
  greater loss.
\end{definition}

\begin{definition}[Erasure monoid and lattice action]
  The set of all erasure events under composition forms a
  \emph{monoid} $(\mathcal{E}, \circ, \varepsilon_\emptyset)$:
  the composition $\varepsilon_s \circ \varepsilon_t$ applies
  $\varepsilon_t$ first, then $\varepsilon_s$;
  the identity is the null erasure $\varepsilon_\emptyset =
  \mathrm{id}$.
  The monoid is not in general a group: erasure events need not
  be invertible.

  The monoid $\mathcal{E}$ acts on the sub-repertoire lattice
  $\mathcal{L}(\Rzero)$ by
  \[
    \mathcal{E} \times \mathcal{L}(\Rzero)
    \;\longrightarrow\; \mathcal{L}(\Rzero),
    \quad
    (\varepsilon, \mathcal{R}) \;\mapsto\; \varepsilon(\mathcal{R})
    \;\subseteq\; \mathcal{R}.
  \]
  An erasure path is an orbit segment of this action.
  Path-dependence of erasure is the statement that two orbits with
  the same endpoints need not traverse the same intermediate
  lattice elements, producing different intermediate admissibility
  geometries.
  The orbit structure of $\mathcal{E}$ on $\mathcal{L}(\Rzero)$
  encodes all possible forgetting trajectories from $\Rzero$.
\end{definition}

\begin{remark}[Structure over elements]
  Meadows~\cite{Meadows2008} argues that in any complex system,
  the interconnections and feedback loops among elements are more
  determinative of system behaviour than the elements themselves.
  The repair-theoretic account instantiates this principle
  precisely: the admissibility geometry $\AO(\mathcal{R})$ and
  the deformation tensor $\Et$ are properties of the
  \emph{structure} of the repair repertoire --- how operators
  compose, what they jointly restore, how their removal cascades
  through the lattice $\mathcal{L}(\Rzero)$ --- not properties
  of individual operators in isolation.
  Erasing one operator changes the geometry not only of the
  distinctions that operator directly maintained but, through
  the causal linkages encoded in the lattice, of every
  distinction whose repair paths passed through it.
  This is the formal analogue of Meadows' observation that
  changing a single connection in a system can alter the
  entire pattern of its behaviour.
\end{remark}

\begin{definition}[Sub-repertoire lattice]
  The set of all sub-repertoires of $\Rzero$, ordered by inclusion,
  forms a complete lattice $(\mathcal{L}(\Rzero), \subseteq)$.
  Meets are intersections; joins are unions (when the union is still
  a valid repertoire).
  An erasure path is a descending chain in this lattice.
\end{definition}

\begin{remark}[Path-dependence of erasure]
  Two erasure paths with the same endpoints
  $\Rzero \to \mathcal{R}_T$ may pass through different intermediate
  repertoires and produce different intermediate admissibility
  geometries, even when $\AO(\mathcal{R}_T)$ is identical.
  This path-dependence is the first indication that erasure is a
  geometric process, not merely a set-theoretic one.
  It is the precursor of the hysteresis phenomenon studied in
  Section~\ref{sec:hysteresis}.
\end{remark}

\begin{definition}[Partial and total repertoire collapse]
  An erasure path reaches \emph{partial collapse} at time $t$ if
  $\Rest(\Rt) \subsetneq \Restzero$ with $\Rest(\Rt) \neq \emptyset$.
  It reaches \emph{total collapse} if $\Rest(\Rt) = \emptyset$.
  Total collapse corresponds to the observer losing all observational
  capacity.
\end{definition}


% ============================================================
\section{Admissibility Contraction}
\label{sec:contraction}
% ============================================================

\begin{theorem}[Admissibility Contraction Theorem]
  Let $\Rt \subseteq \Rzero$ be a repair-repertoire reduction.
  Then
  \[
    \AO(\Rt) \;\subseteq\; \AO(\Rzero).
  \]
  Moreover, if $\Restzt \subsetneq \Restzero$, the inclusion is
  strict and at least one boundary component of $\AO(\Rzero)$
  migrates inward: there exists a point $x \in \AO(\Rzero)$ such
  that $x \notin \AO(\Rt)$.
  \label{thm:contraction}
\end{theorem}

\begin{proof}
  $(\subseteq)$\;
  Suppose $x \in \AO(\Rt)$.
  Then from $x$, every distinction in $\Restzt$ is finitely
  restorable under $\Rt$.
  Since $\Rt \subseteq \Rzero$, every repair path available in
  $\Rt$ is also available in $\Rzero$.
  Hence every distinction in $\Restzt \subseteq \Restzero$ is also
  finitely restorable under $\Rzero$ from $x$, so
  $x \in \AO(\Rzero)$.

  \textbf{Strict inclusion.}\;
  Suppose $d^* \in \Restzero \setminus \Restzt$.
  Then $\CR(d^*, \Rt) = +\infty$: there exists an admissible
  perturbation of $d^*$ from which no finite repair path in $\Rt$
  restores it.
  Let $x_0$ be a state from which $d^*$ was restorable under
  $\Rzero$.
  Since $d^* \notin \Restzt$, there is no finite repair path from
  $x_0$ under $\Rt$ restoring $d^*$.
  Therefore $x_0 \notin \AO(\Rt)$.
  But $x_0 \in \AO(\Rzero)$, establishing strict inclusion.
  By continuity of $\CR(\cdot, \Rzero)$, the set of such points
  $x_0$ is non-empty and open, so a boundary component of
  $\AO(\Rzero)$ has migrated inward.
\end{proof}

\begin{corollary}[Monotonicity of memory under erasure]
  Repair-repertoire reduction monotonically contracts memory:
  \[
    \Rt \subseteq \Rzero
    \quad\Longrightarrow\quad
    \Rest(\Rt) \subseteq \Rest(\Rzero).
  \]
  This is the direct analogue of the Repair Monotonicity Theorem
  of \cite{Flyxion2026a}, now stated in the erasure direction.
\end{corollary}

\begin{corollary}[Boundary migration]
  Under a non-trivial erasure path $(\Rt)_{t \in [0,T]}$,
  the admissibility boundary $\partial\AO(\Rt)$ migrates inward
  monotonically: $\partial\AO(\mathcal{R}_s)$ is strictly interior
  to $\AO(\mathcal{R}_t)$ for $s > t$ whenever the erasure removes
  distinctions.
  The observer's world contracts with every operator removed.
\end{corollary}


% ============================================================
\section{The Erasure Deformation Tensor}
\label{sec:deformation}
% ============================================================

The Admissibility Contraction Theorem is a set-theoretic result.
It tells us that the repair domain shrinks but not how it is
distorted.
To capture distortion we need a tensor.

The key point is that after erasure, the surviving repair domain
$\Restzt$ sits inside the original domain $\Restzero$ via the
inclusion
\[
  \iotat : \Restzt \hookrightarrow \Restzero.
\]
The original metric $\gzero = g_{\Rzero}$ pulls back to
$\iotat^* \gzero$ on $\Restzt$.
The reduced repertoire $\Rt$ induces its own metric $\gt =
g_{\Rt}$ on $\Restzt$.
These need not agree.

\begin{definition}[Erasure deformation tensor]
  The \emph{erasure deformation tensor} at time $t$ is
  \[
    \Et
    \;=\;
    \iotat^* \gzero - \gt,
  \]
  a symmetric $(0,2)$-tensor on $\Restzt$.
  Positive definiteness of $\Et$ at a distinction $d$ means the
  original metric assigns greater distances near $d$ than the
  reduced metric; negative definiteness means the reduced metric is
  locally more spread.
\end{definition}

\begin{remark}
  The pullback construction is essential.
  Writing $\Et = \gzero - \gt$ would compare metrics on different
  spaces, since $\gzero$ is defined on $\Restzero$ and $\gt$ is
  defined on the strictly smaller set $\Restzt$.
  The inclusion $\iotat$ identifies surviving distinctions with
  their images in the original space, so the comparison is
  meaningful: $\Et$ measures how erasure has distorted the geometry
  of what remains, measured against the geometry it formerly had.
\end{remark}

\begin{lemma}[Repair Cost Coercivity]
  Let $d_R^{(t)}$ denote the repair pseudometric induced by the
  repertoire $\Rt$.
  Assume that for every $M < +\infty$, the set of repair paths of
  cost at most $M$ under $\Rt$,
  \[
    \Gamma_M^{(t)}
    \;=\;
    \bigl\{
      \gamma \in \Rt^* \;\big|\; \mathrm{cost}(\gamma) \leq M
    \bigr\},
  \]
  is sequentially compact under the product topology induced by the
  compact-open topology on $\Rt$.
  Suppose $d^* \in \Restzero \setminus \Restzt$ and
  $(d_n)_{n \geq 1}$ is a sequence in $\Restzt$ with
  $d_n \to d^*$ in the $\gzero$-topology.
  Then for every fixed $d_0 \in \Restzt$ not approaching $d^*$,
  \[
    d_R^{(t)}(d_n, d_0) \;\to\; +\infty.
  \]
  \label{lem:coercivity}
\end{lemma}

\begin{proof}
  Suppose for contradiction that the distances remain bounded:
  $d_R^{(t)}(d_n, d_0) \leq M$ for all sufficiently large $n$
  and some $M < +\infty$.
  Then there exists a sequence of repair paths
  $(\gamma_n) \subseteq \Gamma_{M+1}^{(t)}$ with
  $\gamma_n : d_n \leadsto d_0$.

  By the assumed sequential compactness of $\Gamma_{M+1}^{(t)}$,
  a subsequence $(\gamma_{n_k})$ converges to a limit path
  $\gamma^* \in \Gamma_{M+1}^{(t)}$.
  By joint continuity of evaluation in the compact-open topology
  on the locally compact Hausdorff substrate $\Xsub$
  (Axioms 1--3 of \cite{Flyxion2026a}), the terminal states
  converge: $d_{n_k} \to d^*$ implies $\gamma^*$ connects $d^*$
  to $d_0$.
  Since the distinction relation $\Dsub$ is closed
  (Axiom 2 of \cite{Flyxion2026a}) and the terminal state of
  $\gamma_{n_k}$ lies in $\Dsub$, the terminal state of $\gamma^*$
  also lies in $\Dsub$.
  By lower semi-continuity of the cost functional,
  \[
    \mathrm{cost}(\gamma^*) \;\leq\; \liminf_k \mathrm{cost}(\gamma_{n_k})
    \;\leq\; M + 1 \;<\; +\infty.
  \]
  This means $d^*$ is finitely restorable under $\Rt$,
  contradicting $d^* \notin \Restzt$.
  Therefore no uniform bound $M$ exists, and
  $d_R^{(t)}(d_n, d_0) \to +\infty$.
\end{proof}

\begin{remark}[On the path-space compactness hypothesis]
  The hypothesis that $\Gamma_M^{(t)}$ is sequentially compact is
  the path-space analogue of Axiom 3 (compact repair repertoire)
  from \cite{Flyxion2026a}.
  It holds whenever $\Rt$ is a compact set of continuous maps and
  the set of efficient paths of bounded cost can be uniformly
  parameterised --- for instance, when paths have uniformly bounded
  length under the cost density $c$.
  Without this assumption, a bounded-cost sequence of repair paths
  can oscillate indefinitely without converging, leaving the
  contradiction argument incomplete.
  The assumption is satisfied in all concrete applications
  considered in this paper.
\end{remark}

\begin{theorem}[Erasure Deformation Theorem]
  Let $(\Rt)_{t \geq 0}$ be an erasure path satisfying the
  path-space compactness hypothesis of Lemma~\ref{lem:coercivity},
  and let $(d_n)_{n \geq 1}$ be a sequence in $\Restzt$ approaching
  a distinction $d^* \in \Restzero \setminus \Restzt$ in the
  $\gzero$-topology.
  Fix any $d_0 \in \Restzt$ not approaching $d^*$, and let
  $|\cdot|$ denote the operator norm on symmetric $(0,2)$-tensors
  induced by $\gzero$.
  Then
  \[
    |\Et|(d_n) \;\to\; +\infty
    \quad\text{as } d_n \to d^*.
  \]
  Erasure generates geometric singularities at the irrestorability
  boundary.
  \label{thm:deformation}
\end{theorem}

\begin{proof}
  We evaluate the two components of $\Et = \iotat^* \gzero - g(\dRt)$
  along the sequence $(d_n)$, where $g(\dRt)$ denotes the
  $(0,2)$-tensor induced by the repair pseudometric $\dRt =
  d_R^{(t)}$ under the identification of $\Restzt$ as a metric
  space (in the reversible case) or, more generally, the bilinear
  form $g(\dRt)(u,v) = \dRt(u, v) + \dRt(v, u)$ symmetrised from
  the pseudometric.

  \textbf{Pullback component is bounded.}\;
  The pullback metric $(\iotat^* \gzero)(d_n, d_0)$ is the
  $\gzero$-distance between $d_n$ and $d_0$ measured in the
  intact repair space $\Restzero$.
  Since $d_n \to d^*$ in $\Restzero$ and $d_0$ is fixed away from
  $d^*$, the sequence $(\iotat^* \gzero)(d_n, d_0)$ converges to
  $\gzero(d^*, d_0) < +\infty$.
  In particular, $(\iotat^* \gzero)(d_n, d_0)$ is bounded above.

  \textbf{Induced form from reduced pseudometric diverges.}\;
  By Lemma~\ref{lem:coercivity}, $\dRt(d_n, d_0) \to +\infty$.
  Therefore $g(\dRt)(d_n, d_0) \to +\infty$.

  \textbf{Tensor norm diverges.}\;
  The operator norm satisfies
  \[
    |\Et|(d_n)
    \;\geq\;
    \bigl|(\iotat^* \gzero)(d_n, d_0) - g(\dRt)(d_n, d_0)\bigr|.
  \]
  Since $(\iotat^* \gzero)(d_n, d_0)$ is bounded while
  $g(\dRt)(d_n, d_0) \to +\infty$, the right-hand side diverges.
  Therefore $|\Et|(d_n) \to +\infty$.
\end{proof}

\begin{remark}
  Theorem~\ref{thm:deformation} is the geometric content of the
  title.
  Erasure does not merely remove distinctions; it deforms the metric
  structure of surviving distinctions near the lost boundary.
  An observer approaching the irrestorability boundary experiences
  not just loss but increasing distortion of what remains.
  This captures phenomenological features of memory degradation:
  distinctions near the boundary of recall become harder to
  separate, more costly to maintain, and eventually indistinguishable.
\end{remark}

\begin{proposition}[Metric degeneration at the irrestorability boundary]
  Let $\gt$ be a smooth Riemannian metric on $\Restzt$ (in the
  reversible case, where the repair pseudometric satisfies the
  symmetry condition of Proposition~3.5 in \cite{Flyxion2026a}).
  If $|\Et(d_n)| \to +\infty$ along a sequence approaching an
  irrestorability boundary, then either:
  \begin{enumerate}
    \item the smallest eigenvalue of $\gt$ approaches zero
      (metric collapse: the manifold develops a degenerate
      direction), or
    \item the largest eigenvalue of $\gt$ diverges
      (metric blow-up: distances in some direction become
      infinite).
  \end{enumerate}
  In either case, the metric $\gt$ becomes singular at the
  irrestorability boundary.
  \label{prop:metric-degen}
\end{proposition}

\begin{proof}
  Since $|\Et| \to +\infty$ and $\iotat^* \gzero$ is bounded
  (Theorem~\ref{thm:deformation}), the eigenvalues of $\gt$ must
  be unbounded or approach zero along the sequence, for otherwise
  $|\Et|$ would remain bounded.
  The two cases correspond to opposite-signed singularities of the
  metric.
\end{proof}

\begin{corollary}[Forgetting generates metric singularities]
  Under the conditions of Proposition~\ref{prop:metric-degen},
  curvature divergence follows if $\gt$ has bounded derivatives.
  Without this regularity assumption, the weaker statement is that
  the metric becomes singular --- which is a necessary condition for
  curvature blow-up but does not imply it in full generality.
  Forgetting is therefore curvature change in the regular case, and
  metric degeneration in general.
\end{corollary}


% ============================================================
\section{Hysteresis and Irreversibility}
\label{sec:hysteresis}
% ============================================================

The deformation tensor describes the geometry of erasure at a fixed
time.
The hysteresis functional describes the dynamics of restoration
after erasure.

\begin{definition}[Destruction and restoration cost]
  For a distinction $d \in \Restzero$, define:
  \begin{itemize}
    \item the \emph{destruction cost}
      $\CR^{\mathrm{destroy}}(d)$: the minimal cost required to
      remove an operator from $\Rzero$ such that $d$ becomes
      irrestorable;
    \item the \emph{restoration cost}
      $\CR^{\mathrm{restore}}(d, \Rt)$: the minimal cost of
      re-acquiring operators sufficient to return $d$ to
      $\Rest(\Rt)$, starting from the reduced repertoire $\Rt$.
  \end{itemize}
  Both costs admit concrete operational interpretations:
  in biological systems, cost is metabolic energy or cellular
  resources ($\CR^{\mathrm{destroy}}$ = energy to eliminate a
  synaptic pathway; $\CR^{\mathrm{restore}}$ = energy to grow a
  new one); in institutional systems, cost is organisational
  resources ($\CR^{\mathrm{destroy}}$ = cost of decommissioning
  an archive or retiring a specialist;
  $\CR^{\mathrm{restore}}$ = cost of rebuilding equivalent
  capacity from scratch); in linguistic systems, cost is
  measurable in speaker-years
  ($\CR^{\mathrm{destroy}}$ = attrition required to make a
  distinction irrestorable; $\CR^{\mathrm{restore}}$ =
  revitalisation effort required to restore it).
  The key asymmetry is that destruction cost is measured going
  down the sub-repertoire lattice, while restoration cost is
  measured going up from $\Rt$; the hysteresis functional
  captures this asymmetry in a single signed quantity.
\end{definition}

\begin{definition}[Hysteresis functional]
  The \emph{hysteresis functional} of distinction $d$ under
  erasure path $(\Rt)$ is
  \[
    \Hyster_t(d)
    \;=\;
    \CR^{\mathrm{restore}}(d, \Rt)
    - \CR^{\mathrm{destroy}}(d).
  \]
  \label{def:hysteresis}
\end{definition}

Three regimes characterise the hysteresis functional:

\begin{itemize}
  \item $\Hyster_t(d) = 0$: \emph{reversible erasure}.
    Restoration costs exactly as much as destruction.
    The distinction is forgotten and recovered with equal effort.
    This regime governs ordinary perturbation within the
    admissibility interior.

  \item $0 < \Hyster_t(d) < +\infty$: \emph{partial hysteresis}.
    Restoration is more costly than destruction.
    This is the generic biological and institutional regime:
    synaptic weakening is reversible but synaptic elimination
    requires costly re-growth; institutional procedures once
    suspended require substantial effort to reinstate.

  \item $\Hyster_t(d) = +\infty$: \emph{permanent forgetting}.
    The distinction cannot be restored from $\Rt$ at any finite
    cost.
\end{itemize}

\begin{theorem}[Permanent Forgetting Theorem]
  $\Hyster_t(d) = +\infty$ if and only if
  $d \notin \Rest(\Rt)$.
  Infinite hysteresis is equivalent to the distinction having
  crossed the irrestorability boundary.
  \label{thm:permanent}
\end{theorem}

\begin{proof}
  $(\Rightarrow)$\;
  If $\Hyster_t(d) = +\infty$, then
  $\CR^{\mathrm{restore}}(d, \Rt) = +\infty$: no finite
  re-acquisition of operators restores $d$ from $\Rt$.
  This means $d \notin \Rest(\Rt)$.

  $(\Leftarrow)$\;
  If $d \notin \Rest(\Rt)$, then $d$ is not finitely restorable
  under $\Rt$.
  Any sequence of operator acquisitions that would restore $d$
  must ultimately recover a repertoire $\mathcal{R}'$ with
  $d \in \Rest(\mathcal{R}')$.
  The cost of doing so from $\Rt$ is $\CR^{\mathrm{restore}}(d,
  \Rt) = +\infty$ if no finite extension of $\Rt$ suffices ---
  which holds whenever the operators needed to restore $d$ are
  themselves irrecoverable (as in biological cell death or archive
  destruction).
  Hence $\Hyster_t(d) = +\infty$.
\end{proof}

\begin{remark}
  The Permanent Forgetting Theorem connects the dynamical notion
  (infinite restoration cost) to the geometric notion (crossing the
  admissibility boundary).
  It is the hysteresis analogue of the Repair Boundary Theorem of
  \cite{Flyxion2026a}: where that theorem characterised the boundary
  via diverging repair cost, this theorem characterises it via
  diverging restoration cost after the boundary has been crossed.
\end{remark}

\begin{proposition}[Monotonicity of hysteresis under progressive erasure]
  Along an erasure path $(\Rt)_{t \geq 0}$, the hysteresis
  functional is non-decreasing in $t$:
  \[
    s \geq t \quad\Longrightarrow\quad
    \Hyster_s(d) \geq \Hyster_t(d)
  \]
  for all $d \in \Rest(\Rzero)$.
  Further erasure never makes restoration easier.
\end{proposition}

\begin{proof}
  Since $\mathcal{R}_s \subseteq \Rt$ for $s \geq t$, the set of
  operators available for restoration at time $s$ is a subset of
  those available at time $t$.
  Therefore $\CR^{\mathrm{restore}}(d, \mathcal{R}_s) \geq
  \CR^{\mathrm{restore}}(d, \Rt)$, giving
  $\Hyster_s(d) \geq \Hyster_t(d)$.
\end{proof}

\begin{theorem}[Hysteresis--Boundary Correspondence]
  Let $d \in \Restzero$.
  Then
  \[
    \Hyster_t(d) = +\infty
    \quad\Longleftrightarrow\quad
    \inf_{\gamma}
    \Bigl(
      \sup_{x \in \gamma} \CR(x,\, \Rt)
    \Bigr)
    = +\infty,
  \]
  where the infimum is taken over all restoration trajectories
  $\gamma$ from $\Rt$ to any repertoire $\mathcal{R}' \supseteq \Rt$
  with $d \in \Rest(\mathcal{R}')$.
  Permanent forgetting is equivalent to every restoration trajectory
  encountering a minimax barrier of infinite repair cost.
  \label{thm:hysteresis-boundary}
\end{theorem}

\begin{proof}
  $(\Rightarrow)$\;
  Suppose $\Hyster_t(d) = +\infty$, so
  $\CR^{\mathrm{restore}}(d, \Rt) = +\infty$.
  Suppose for contradiction that there exists a restoration
  trajectory $\gamma^*$ with
  \[
    \sup_{x \in \gamma^*} \CR(x,\, \Rt) \;<\; +\infty.
  \]
  Then every point along $\gamma^*$ lies at finite repair cost, and
  the total restoration cost accumulated along $\gamma^*$ is bounded
  above by the path length times this finite supremum.
  Therefore $\CR^{\mathrm{restore}}(d, \Rt) < +\infty$, contradicting
  the hypothesis.
  Hence no such trajectory exists: every restoration trajectory
  $\gamma$ satisfies $\sup_{x \in \gamma} \CR(x, \Rt) = +\infty$.
  Taking the infimum over all trajectories gives
  $\inf_\gamma \sup_{x \in \gamma} \CR(x, \Rt) = +\infty$.

  $(\Leftarrow)$\;
  Suppose $\inf_\gamma \sup_{x \in \gamma} \CR(x, \Rt) = +\infty$.
  Then for every restoration trajectory $\gamma$,
  $\sup_{x \in \gamma} \CR(x, \Rt) = +\infty$.
  By the Repair Boundary Theorem of \cite{Flyxion2026a}, points
  where $\CR(\cdot, \Rt) \to +\infty$ are precisely points on the
  irrestorability boundary $\partial\AO(\Rt)$.
  Every trajectory therefore accumulates at boundary points of
  infinite repair cost and cannot reach $d$ at finite total cost.
  Hence $\CR^{\mathrm{restore}}(d, \Rt) = +\infty$, giving
  $\Hyster_t(d) = +\infty$.
\end{proof}

\begin{remark}
  The Hysteresis--Boundary Correspondence gives permanent forgetting
  a variational characterisation: $\Hyster_t(d) = +\infty$ if and
  only if the minimax repair cost over all restoration trajectories
  is infinite.
  This is the exact analogue of barrier-crossing formulations in
  control theory, large deviations theory, and energy landscape
  analysis, where a transition between states is impossible if the
  infimum of the worst-case cost over all paths is infinite.
  The theorem unifies admissibility geometry (the boundary),
  dynamical cost (the hysteresis functional), and the Permanent
  Forgetting Theorem into a single minimax equivalence.
  Every concept introduced in this paper --- contraction,
  deformation, hysteresis, permanent forgetting --- is now
  a consequence of a single geometric fact: that irrestorability
  boundaries are minimax barriers of infinite repair cost.
\end{remark}

\begin{theorem}[Boundary Migration]
  Let $(\Rt)_{t \geq 0}$ be a non-trivial erasure path.
  For each $t > 0$ such that $\Restzt \subsetneq \Restzero$,
  there exists an open set $U_t \subset \Azero$ with
  $U_t \cap \At = \emptyset$ and $U_t \cap \partial\At \neq
  \emptyset$: the irrestorability boundary $\partial\At$ has
  migrated inward into the former interior of $\Azero$.
  \label{thm:migration}
\end{theorem}

\begin{proof}
  By the Admissibility Contraction Theorem, $\At \subseteq \Azero$
  with strict inclusion when $\Restzt \subsetneq \Restzero$.
  Let $x_0 \in \Azero \setminus \At$.
  Since $\Azero$ is the closure of its interior, points arbitrarily
  close to $x_0$ were formerly in the interior of $\Azero$.
  As $\At \not\ni x_0$, the boundary $\partial\At$ lies between
  $x_0$ and the interior of $\At$, establishing $U_t$.
\end{proof}

\begin{definition}[Boundary migration measure]
  The \emph{boundary migration measure} at time $t$ is
  \[
    V(t)
    \;=\;
    \inf_{x \in \partial\At}\;
    \inf_{y \in \partial\Azero}\;
    d_R^{(0)}(x, y),
  \]
  where $d_R^{(0)}$ is the original repair pseudometric and the
  double infimum gives the point-to-set distance from the current
  irrestorability boundary $\partial\At$ to the original boundary
  $\partial\Azero$ in the original repair geometry.
  $V(t)$ is non-negative with $V(t) > 0$ whenever the boundary has
  migrated inward.
  Whether $V(t)$ is differentiable in $t$, and what the derivative
  would represent as a migration speed, depends on regularity
  assumptions about the boundary evolution not established here;
  this is an open problem.
\end{definition}

\begin{remark}
  Theorem~\ref{thm:migration} reformulates forgetting as a
  moving-front problem.
  The irrestorability boundary advances inward through the
  admissibility manifold, sweeping formerly admissible states into
  the irrestorable region.
  The migration measure $V(t)$ quantifies how far the front has
  advanced without requiring smoothness of the boundary, making it
  applicable even when the boundary fragments or develops
  topological complexity under erasure.
  Biological analogues include the penumbra in stroke and the
  retreat of phonological distinctions as speaker populations
  decline.
\end{remark}


% ============================================================
\section{Topological Memory}
\label{sec:topology}
% ============================================================

The deformation tensor and hysteresis functional are local objects:
they describe geometry near individual distinctions and boundaries.
To capture global structure of memory loss, we need topology.

The metric and topological levels of the framework answer different
questions and should be clearly distinguished.
The metric structure --- the repair pseudometric $\dRt$ and the
deformation tensor $\Et$ --- answers the question: \emph{how costly
is access to a distinction?}
It is quantitative, local, and sensitive to the specific distances
between distinctions in the repair domain.
The topological structure --- the homology groups $H_k(\At)$ and
the persistence module $PH_*(\AO)$ --- answers the question:
\emph{which distinctions remain mutually accessible at all?}
It is qualitative, global, and invariant under homeomorphism.
These are independent.
Large metric deformation can occur without any topological change
(if the repair domain contracts but remains connected and simply
connected).
Conversely, topological change --- fragmentation, hole creation ---
can occur with bounded local metric deformation.
The full geometry of forgetting requires both levels: topology
describes the qualitative structure of what remains, metric geometry
describes the quantitative structure of how easily it is accessed.

The erasure path $(\Rt)_{t \geq 0}$ generates a filtration of
admissibility manifolds:
\[
  \Azero \;\supseteq\; \AO(\mathcal{R}_1) \;\supseteq\; \AO(\mathcal{R}_2)
  \;\supseteq\; \cdots
\]
This filtration is the natural input to persistent homology.
The individual homology groups $H_k(\At)$ record topology at a
single time $t$; the persistence module records how that topology
evolves across the entire erasure path.

\begin{definition}[Persistence module of memory]
  The \emph{persistence module} of the observer's memory under the
  erasure path $(\Rt)$ is the collection of maps
  \[
    \mathbf{PH}_k
    \;=\;
    \bigl(H_k(\At;\, \mathbb{Z}),\;
    \phi_{s,t} : H_k(\At) \to H_k(\AO(\mathcal{R}_s))\bigr)_{s \geq t},
  \]
  where $\phi_{s,t}$ is induced by the inclusion $\AO(\mathcal{R}_s)
  \hookrightarrow \At$.
  The \emph{persistent topological memory} is
  \[
    PH_*(\AO)
    \;=\;
    \bigoplus_{k \geq 0} \mathbf{PH}_k.
  \]
\end{definition}

The persistence diagram of $PH_*(\AO)$ records, for each
homological feature, its birth time $t_b$ (the erasure time at
which it first appears) and death time $t_d$ (the time at which
it is destroyed by further erasure).
Features with large lifespan $t_d - t_b$ represent structurally
robust repair capacity; features with small lifespan represent
fragile repair cycles sensitive to operator loss.

\begin{remark}[On homology coefficients]
  We use $\mathbb{Z}$ coefficients throughout.
  With $\mathbb{Z}$ coefficients, $H_0(\At; \mathbb{Z})$ is the
  free abelian group on connected components of $\At$, so
  $\mathrm{rank}(H_0) = $ number of components --- exactly the
  quantity relevant for the Fragmentation Theorem.
  With $\mathbb{Z}/2$ coefficients the rank would agree for
  orientable manifolds, but $\At$ need not be orientable after
  erasure.
  Integer coefficients are the natural choice when we wish to
  count components and loops without orientability assumptions,
  which is the appropriate setting for admissibility manifolds
  whose topology may be altered by non-orientable erasure events.
\end{remark}

\begin{definition}[Topological memory]
  The \emph{topological memory} of an observer at time $t$ is
  the graded object
  \[
    \Mem_t
    \;=\;
    \bigoplus_{k \geq 0} H_k(\At;\, \mathbb{Z}),
  \]
  where $\At = \AO(\Rt)$ is the admissibility manifold at time $t$.
  The full memory across the erasure path is $PH_*(\AO)$.
\end{definition}

Each homology group captures a distinct mode of persistent memory
structure:

\begin{itemize}
  \item $H_0(\At)$: connected components of the admissibility
    manifold.
    A single component corresponds to a globally connected repair
    space; multiple components indicate \emph{repair domain
    fragmentation} --- regions of the distinction space that can
    no longer be jointly maintained.

  \item $H_1(\At)$: loops in the admissibility manifold.
    Non-trivial loops correspond to \emph{repair cycles} ---
    closed sequences of repair operations that return to their
    starting point.
    Their birth and death in $PH_1$ records the creation and
    destruction of circular repair redundancy under erasure.

  \item $H_k(\At)$ for $k \geq 2$: higher-order topological
    features corresponding to repair structures with more complex
    boundary behaviour.
\end{itemize}

\begin{definition}[Topological forgetting]
  A mode of \emph{topological forgetting} occurs when an erasure
  event changes the homology of $\At$:
  \begin{itemize}
    \item \emph{Fragmentation}: $\mathrm{rank}(H_0(\At))$ increases.
    \item \emph{Hole creation}: $\mathrm{rank}(H_1(\At))$ increases.
    \item \emph{Homology collapse}: $\mathrm{rank}(H_k(\At))$
      decreases for some $k$.
  \end{itemize}
  All three modes are visible in the persistence diagram as
  births, long-lived features, and deaths respectively.
\end{definition}

\begin{theorem}[Fragmentation Theorem]
  Let $(\Rt)_{t \geq 0}$ be an erasure path that removes a
  distinction $d^* \in \Restzero$.
  If $d^*$ was a cut-point of $\Azero$, then the erasure event
  that removes $d^*$ increases $\mathrm{rank}(H_0(\At))$ by at
  least one.
  Fragmentation is the topological signature of the loss of a
  structurally critical distinction.
  \label{thm:homology-change}
\end{theorem}

\begin{proof}
  If $d^*$ is a cut-point of $\Azero$, then $\Azero \setminus
  \{d^*\}$ is disconnected.
  After the erasure event making $d^*$ irrestorable, $\At$ is
  contained in $\Azero \setminus \{d^*\}$ and is therefore
  disconnected, increasing $\mathrm{rank}(H_0(\At))$.
\end{proof}

\begin{theorem}[Hole Creation Theorem]
  Let $U \subset \At$ be a contractible region of the admissibility
  manifold whose repair operators are entirely removed by an erasure
  event, so that $U \cap \AO(\mathcal{R}_{t+\varepsilon}) =
  \emptyset$ for all small $\varepsilon > 0$.
  If the removal of $U$ does not disconnect $\At$ but creates a
  non-contractible loop in $\At \setminus U$, then
  \[
    \mathrm{rank}(H_1(\AO(\mathcal{R}_{t+\varepsilon})))
    \;>\;
    \mathrm{rank}(H_1(\At)).
  \]
  Erasure can increase topological complexity even while decreasing
  admissible volume.
  \label{thm:hole-creation}
\end{theorem}

\begin{proof}
  The removed region $U$ becomes inaccessible under the reduced
  repertoire.
  The complement $\At \setminus U$ contains a loop surrounding $U$
  that is not the boundary of any admissible $2$-chain in
  $\At \setminus U$: no $2$-chain can fill the loop because $U$ is
  no longer part of the admissible space.
  Hence a new non-trivial homology class appears in $H_1$, and the
  rank increases.
\end{proof}

\begin{remark}
  Theorem~\ref{thm:hole-creation} captures a counter-intuitive
  phenomenon: an observer can become topologically more complex
  through forgetting.
  The emergence of inaccessible repair cycles --- holes in the
  admissibility manifold --- represents a kind of structural
  complication that progressive erasure introduces even as the
  total admissible volume shrinks.
  This is the geometric analogue of the observation that partial
  memory loss can produce confabulatory behaviour: the system
  develops topological constraints it cannot see and cannot
  navigate around.
\end{remark}

\begin{remark}
  Ray~Barman et al.~\cite{RayBarman2026} observe a sharp phase
  transition in $H_0$ and a peak in $H_1$ of semantic embedding
  spaces at a critical connectivity threshold, with $H_1$ reaching
  $534 \pm 24$ at the transition.
  This is an empirical instance of the Hole Creation Theorem:
  at the critical threshold, the removal of connectivity between
  semantic clusters creates non-contractible loops in the surviving
  space, exactly as Theorem~\ref{thm:hole-creation} predicts.
  The persistence diagram of their embedding space is the empirical
  correlate of $PH_*(\AO)$ for a retrieval observer undergoing
  progressive interference-driven erasure.
\end{remark}

\begin{definition}[Persistent topological memory]
  A topological feature of $PH_*(\AO)$ \emph{persists} if its
  lifespan $t_d - t_b$ in the persistence diagram exceeds a
  threshold $\tau > 0$.
  Persistent features are the structurally robust components of
  memory; ephemeral features are fragile repair cycles that erasure
  quickly destroys.
  The persistence diagram provides a complete summary of which
  repair structures survive progressive operator loss.
\end{definition}


% ============================================================
\section{Biological Forgetting}
\label{sec:bio}
% ============================================================

The biological substrate for distinction-maintaining repair is
the neural system.
Repair operators include synaptic strengthening, protein synthesis,
glial repair, and the broader repertoire of cellular maintenance
processes.
Erasure events include synaptic weakening, synaptic elimination,
cell death, and axonal degeneration.

\begin{example}[Synaptic decay and partial hysteresis]
  Long-term potentiation (LTP) is a repair operator that
  strengthens synaptic connections, increasing the probability of
  successful distinction restoration.
  Synaptic weakening below the LTP threshold is reversible:
  $0 < \Hyster_t(d) < +\infty$.
  Synaptic elimination by pruning or excitotoxicity is in general
  irreversible: $\Hyster_t(d) = +\infty$ if the cellular
  machinery for synapse formation is also lost.
  The transition between these regimes is a biological instance of
  the boundary crossing described by Theorem~\ref{thm:permanent}.
\end{example}

\begin{example}[Neural injury and admissibility contraction]
  Focal cortical lesions remove repair operators acting on specific
  regions of the substrate.
  By the Admissibility Contraction Theorem, the admissibility
  manifold of the injured system is a strict subset of the intact
  system's manifold.
  The deformation tensor $\Et$ is non-trivial near the lesion
  boundary: distinctions that were easily maintained before the
  lesion become increasingly costly as they approach the injured
  region.
  This is the geometric correlate of the penumbra effect in
  stroke: the region surrounding the infarct in which repair is
  expensive but not yet impossible.
\end{example}

\begin{example}[Developmental pruning and homology change]
  Synaptic pruning during development eliminates a substantial
  fraction of early repair operators.
  If pruned synapses were cut-points of the admissibility manifold,
  Theorem~\ref{thm:homology-change} predicts fragmentation of the
  repair domain.
  Conversely, if pruning eliminates redundant paths while
  preserving connectivity, $H_0(\At)$ is unchanged and only
  higher-order homology is affected.
  The distinction between pathological and healthy pruning may
  correspond precisely to whether pruning changes $H_0$ (harmful)
  or only higher homology (reorganising without fragmenting).
\end{example}


% ============================================================
\section{Institutional Forgetting}
\label{sec:institutional}
% ============================================================

Institutions are observers whose repair repertoires consist of
codified procedures, trained personnel, and maintained
infrastructure.
Erasure events include staff attrition, procedural reform,
budget cuts, and physical destruction of facilities.

\begin{example}[Archive destruction]
  A physical archive is a repair operator: it restores distinctions
  between historical states by maintaining retrieval paths from
  present queries to past records.
  Destruction of an archive removes those operators.
  By the Permanent Forgetting Theorem, distinctions whose only
  repair paths passed through the destroyed archive acquire
  $\Hyster_t(d) = +\infty$: they cannot be restored at any finite
  cost from the reduced repertoire.
  The distinctions do not become false; they become unrestorable.
  This is the precise sense in which archive destruction is
  epistemically irreversible.
\end{example}

\begin{example}[Language death]
  A living language is a repair repertoire for semantic and
  phonological distinctions.
  Language death is an erasure path: as speakers die or shift to
  other languages, the repair operators (speakers capable of
  producing and recognising the language's distinctions) are
  progressively removed.
  The deformation tensor $\Et$ is non-trivial near the distinctions
  specific to the dying language: they become increasingly costly
  to maintain as the speaker population falls.
  Total collapse of the speaker population corresponds to
  $\Restzt = \emptyset$: topological total forgetting.
  Language revitalisation is the re-acquisition of repair operators,
  with a hysteresis cost determined by how far along the erasure
  path the language has progressed.
\end{example}

\begin{example}[Institutional collapse]
  The collapse of a legal or governmental institution removes
  procedural repair operators.
  The Paradigm Shift Theorem of \cite{Flyxion2026a} showed that
  non-trivial repertoire change produces incomparable restorable
  sets.
  Institutional collapse is the limiting case: the entire
  repertoire is lost, and the new institution that replaces it
  typically has a different restorable set, not a subset.
  Reconstruction of collapsed institutions is path-dependent:
  different reconstruction paths through the sub-repertoire lattice
  produce different admissibility geometries even when the
  nominal repertoire is nominally restored.
  Meadows~\cite{Meadows2008} identifies \emph{eroding goals} and
  \emph{policy resistance} as characteristic system traps: when
  a system loses the feedback loops that maintain a standard, the
  standard itself erodes rather than the gap between standard and
  performance being corrected.
  In repair-theoretic terms, eroding goals correspond to a slow
  erasure path in which the operators that enforce a distinction
  are progressively weakened before being eliminated, with the
  distinction boundary migrating inward --- precisely the
  moving-front dynamics of Theorem~\ref{thm:migration} --- before
  the distinction becomes permanently irrestorable.
\end{example}


% ============================================================
\section{Scientific Forgetting}
\label{sec:science}
% ============================================================

Scientific knowledge is maintained by a collective repair
repertoire: experimental techniques, trained practitioners,
instrumentation, and institutional infrastructure.
Scientific forgetting is the loss of repair operators from this
collective repertoire.

\begin{example}[Loss of experimental technique]
  Certain experimental techniques become inaccessible when the
  practitioners who developed them retire without training
  successors.
  The technique is an operator $r \in \Rzero$; its loss is an
  erasure event.
  Distinctions whose restoration depended exclusively on that
  technique acquire $\Hyster_t(d) = +\infty$ if the technique
  cannot be reconstructed from surviving operators.
  This is institutionally invisible: the distinctions remain
  nominally part of the scientific record but are practically
  unrestorable.
\end{example}

\begin{example}[Paradigm transitions as structured erasure]
  The Paradigm Shift Theorem of \cite{Flyxion2026a} characterised
  paradigm transitions as producing incomparable restorable sets.
  Under the erasure framework, a paradigm transition is an erasure
  path followed by a repertoire acquisition path: old operators
  are removed, new ones are added.
  The net effect on the deformation tensor depends on whether the
  new operators can restore the distinctions made irrestorable by
  the old operators' removal.
  If not, the transition involves permanent scientific forgetting
  --- the Kuhnian claim that paradigm shifts make some distinctions
  permanently invisible --- formalised as $\Hyster_T(d) = +\infty$
  for distinctions in $\Restzero \setminus \Rest(\mathcal{R}_T)$.
\end{example}

\begin{example}[Replication crisis as repair failure]
  The replication crisis in empirical sciences can be recast as
  a collective repair failure: experimental results that were
  nominally in $\Restzero$ are found, upon attempted restoration
  under independent operators, to have $\CR(d, \Rzero) = +\infty$
  --- they were never genuinely restorable, only apparently so.
  The crisis is the discovery that the admissibility boundary was
  mislocated: distinctions believed to be in the interior of
  $\Azero$ were in fact on or outside it.
  %% EXPAND: connection to distinction-reorganisation papers.
\end{example}


% ============================================================
\section{Consequences, Open Problems, and Future Directions}
\label{sec:consequences}
% ============================================================

The framework developed here yields a unified account of phenomena
that are conventionally treated as belonging to different domains.

\medskip\noindent
\textbf{Memory is geometry.}\;
The metrised repair space $(\Rest(\mathcal{R}), d_R)$ is the
formal object corresponding to what an observer remembers.
Memory loss is not subtraction of stored items but deformation of
this metric space, measured by the tensor $\Et$.

\medskip\noindent
\textbf{The right objects are invariants, not components.}\;
The deformation tensor $\Et$ is a symmetric $(0,2)$-tensor; its
components in any particular basis carry no direct physical
meaning.
The meaningful objects are invariants: the operator norm $|\Et|$
measures total deformation; the eigenvalues of $\Et$ (relative to
$\gzero$) identify principal directions of forgetting; the trace
$\mathrm{tr}(\Et)$ measures mean distortion; the determinant
measures volume change in the surviving repair domain.
This is precisely how strain tensors are interpreted in continuum
mechanics: not component by component, but through invariant
quantities such as principal strains, volumetric strain, and
deviatoric strain.
A fading memory corresponds not to a particular tensor component
becoming large but to a specific invariant --- likely $|\Et|$ or
a principal eigenvalue --- diverging in the direction associated
with the distinction being lost.

\medskip\noindent
\textbf{Forgetting is irreversible in the generic case.}\;
The hysteresis functional $\Hyster_t(d)$ is non-decreasing along
any erasure path.
For biological and institutional observers, the operators most
likely to be lost --- cell death, archive destruction, speaker
extinction --- are precisely those for which
$\Hyster_t(d) = +\infty$.
Reversible forgetting is the special case, not the default.

\medskip\noindent
\textbf{Restoration is path-dependent.}\;
Two observers that have undergone different erasure paths to the
same reduced repertoire may have different restoration costs,
because intermediate states of the admissibility manifold affect
what residual structure is available for reconstruction.
This connects the framework to the broader distinction-history
program: what matters is not the current state but the trajectory.

\medskip\noindent
\textbf{Forgetting has two independent modes.}\;
Metric deformation ($|\Et| \to \infty$) and topological change
($H_*(\At)$ changing) are independent phenomena.
The full geometry of forgetting requires both: metric deformation
describes how costly access becomes; topological change describes
which distinctions lose mutual accessibility entirely.
A storage theory of memory can represent neither.

\subsection*{Open Problems}

\begin{enumerate}

  \item \textbf{Boundary migration speed.}\;
    The migration measure $V(t) = \inf_{x \in \partial\At}
    \inf_{y \in \partial\Azero} d_R^{(0)}(x, y)$
    is defined but its time
    derivative is not established.
    Under what regularity conditions on the erasure path and the
    admissibility boundary is $V(t)$ differentiable?
    When it is, $\dot{V}(t)$ is the migration speed --- a scalar
    quantifying how fast the irrestorability front advances.
    Is there an equation of motion for the front analogous to the
    eikonal equation in geometric optics?

  \item \textbf{Rate of deformation.}\;
    How does $\Et$ evolve along an erasure path?
    Is there a flow equation $\partial_t \Et = F(\Et, g_0, \Rt)$?
    The analogy with Ricci flow is suggestive: curvature-driven
    deformation might model how metric singularities propagate
    through the repair domain after a localised erasure event.
    Ricci flow develops singularities in finite time, which would
    correspond to the onset of permanent forgetting.

  \item \textbf{Critical repair operators.}\;
    The Fragmentation Theorem identifies cut-points as structurally
    critical distinctions; the Hole Creation Theorem identifies
    contractible removable regions.
    Characterise the full set of \emph{critical repair operators}:
    those whose removal changes $H_*(\At)$.
    Are there operators that stabilise the manifold topology
    against the removal of other operators?
    If so, they are candidates for targeted preservation in
    biological, archival, and linguistic contexts.

  \item \textbf{Minimal restoration.}\;
    Given a reduced repertoire $\Rt$ and a target distinction set
    $D \subseteq \Restzero \setminus \Restzt$, what is the minimal
    operator set that must be re-acquired to restore $D$?
    This is an inverse problem on the sub-repertoire lattice
    $\mathcal{L}(\Rzero)$.
    The orbit structure of the erasure monoid
    (Section~\ref{sec:erasure}) constrains the answer: not all
    paths from $\Rt$ to a repertoire restoring $D$ have equal cost.

  \item \textbf{Erasure and the interference matrix.}\;
    How does the interference matrix $\IntMat$ of
    \cite{Flyxion2026a} evolve under erasure?
    Does operator loss preferentially increase positive eigenvalues
    (more repair conflict among surviving distinctions) or decrease
    negative eigenvalues (less synergy)?
    The effective dimensionality concentration found by
    Ray~Barman et al.~\cite{RayBarman2026} --- all tested embedding
    models concentrating variance into roughly 16 effective
    dimensions regardless of nominal dimensionality --- suggests
    that biological and artificial memory systems operate in a regime
    where a small number of operators account for most repair
    capacity, making them catastrophically vulnerable to targeted
    erasure.
    A prediction: erasure of the operators corresponding to the
    top eigenvectors of $\IntMat$ should produce disproportionate
    fragmentation relative to erasure of random operators.

  \item \textbf{Optimal erasure paths.}\;
    On $\mathcal{L}(\Rzero)$, characterise the erasure paths that
    minimise total deformation $\int_0^T |\Et|\, \mathrm{d}t$
    subject to a fixed endpoint $\mathcal{R}_T$.
    Such paths represent maximally graceful forgetting.
    The dual problem --- paths that maximise $\int_0^T |\Et|\,
    \mathrm{d}t$ --- characterises maximally disruptive forgetting,
    relevant to adversarial institutional erasure.

\end{enumerate}

\subsection*{Future Directions}

The following questions lie beyond the scope of this paper but
are natural extensions of its framework.

\medskip\noindent
\textbf{Stratified admissibility manifolds.}\;
The current framework assumes $\At$ is a topological manifold.
The Fragmentation and Hole Creation Theorems predict, however,
that progressive erasure can produce non-manifold spaces:
boundaries may become fractal, strata may appear at different
dimensions, and the singular set of $\At$ may itself have
non-trivial topology.
The natural generalisation is to work in the category of
\emph{stratified spaces}~\cite{Munkres2000}, where the metric
tensor lives on each stratum and the homology groups are computed
globally.
This would accommodate the full range of singularities that
erasure generates without requiring smoothness assumptions that
cannot be established in general.

\medskip\noindent
\textbf{Thermodynamics of erasure.}\;
The framework uses repair cost as a primitive but does not
specify its physical nature.
Landauer's principle~\cite{Landauer1961} establishes that the
erasure of one bit of information requires at least $kT \ln 2$
of work.
If distinctions are identified with bits and erasure events are
identified with logical irreversibility, the hysteresis functional
$\Hyster_t(d)$ should satisfy a Landauer-type lower bound:
$\CR^{\mathrm{destroy}}(d) \geq kT \ln 2$ for any distinction
whose erasure is informationally irreversible.
Whether the deformation tensor $\Et$ is conjugate to a
stress-energy tensor in the physical sense --- whether there is
a genuine variational principle underlying erasure deformation
--- is an open question connecting this framework to the
thermodynamics of computation and the physics of information.

\medskip\noindent
\textbf{Continuous operator weights.}\;
The current framework treats repair operators as present or
absent: $r \in \mathcal{R}$ or $r \notin \mathcal{R}$.
In biological systems, synaptic connections are weakened
continuously before elimination; in institutional systems,
procedures are degraded incrementally before abandonment.
A natural extension defines operator weights $w_r \in [0, 1]$
and replaces the discrete lattice $\mathcal{L}(\Rzero)$ with a
continuous space of weighted repertoires.
The deformation tensor $\Et$ would then vary continuously with
the weights, and its singularities would correspond to weight
thresholds at which distinctions cross from restorable to
irrestorable.
Partial hysteresis ($0 < \Hyster_t(d) < +\infty$) is precisely
the regime this extension would formalise rigorously.

\medskip\noindent
\textbf{Observer identity deformation.}\;
The current paper assumes observer persistence: the observer
before and after the erasure path is the same observer with
reduced capacity.
For sufficiently severe forgetting, this assumption fails.
A stroke patient with extensive cortical damage, a language
community reduced to one elderly speaker, or an institution that
has lost all organisational memory may constitute a different
observer rather than a depleted version of the same one.
The natural question is whether observer identity is itself
definable as a path-connected component of the category of
admissibility manifolds: two observers are the same if there
exists an erasure path connecting their repertoires within a
single connected component; they are different if no such path
exists.
This would make observer identity a topological invariant of the
admissibility geometry, and observer deformation --- the passage
between observer-identity classes --- a phenomenon requiring the
full categorical machinery to describe.
It is probably the deepest question to which the current
framework points.

\medskip
Erasure is not information loss.
It is the deformation of a world.
What is forgotten is not gone; it has moved to a boundary where
repair costs diverge, where the metric of what remains is singular,
and where the topology of the admissible has been irreversibly
changed.
The observer that forgets is not smaller.
It is differently shaped.

% ============================================================
\begin{thebibliography}{99}

\bibitem{Flyxion2026a}
Flyxion.
\newblock Observability is restorability: Restorable distinctions and the
  geometry of perceptual access.
\newblock Preprint, 2026.

\bibitem{RayBarman2026}
Ray~Barman, Sambartha, Andrey Starenky, Sophia Bodnar, Nikhil Narasimhan,
  and Ashwin Gopinath.
\newblock The geometry of forgetting.
\newblock \emph{arXiv}:2604.06222, 2026.

\bibitem{Ashby1956}
Ashby, W.~Ross.
\newblock \emph{An Introduction to Cybernetics}.
\newblock Chapman \& Hall, London, 1956.

\bibitem{Kuhn1962}
Kuhn, Thomas~S.
\newblock \emph{The Structure of Scientific Revolutions}.
\newblock University of Chicago Press, Chicago, 1962.

\bibitem{Munkres2000}
Munkres, James~R.
\newblock \emph{Topology}.
\newblock Prentice Hall, Upper Saddle River, NJ, 2nd edition, 2000.

\bibitem{RockafellarWets1998}
Rockafellar, R.~Tyrrell, and Roger~J.-B. Wets.
\newblock \emph{Variational Analysis}.
\newblock Springer, Berlin, 1998.

\bibitem{Villani2009}
Villani, C\'edric.
\newblock \emph{Optimal Transport: Old and New}.
\newblock Springer, Berlin, 2009.

\bibitem{Hatcher2002}
Hatcher, Allen.
\newblock \emph{Algebraic Topology}.
\newblock Cambridge University Press, Cambridge, 2002.

\bibitem{EdelsbrunnerHarer2010}
Edelsbrunner, Herbert, and John~L. Harer.
\newblock \emph{Computational Topology: An Introduction}.
\newblock American Mathematical Society, Providence, RI, 2010.

\bibitem{Sontag1998}
Sontag, Eduardo~D.
\newblock \emph{Mathematical Control Theory}.
\newblock Springer, New York, 1998.

\bibitem{Landauer1961}
Landauer, Rolf.
\newblock Irreversibility and heat generation in the computing process.
\newblock \emph{IBM Journal of Research and Development}, 5(3):183--191, 1961.

\bibitem{Meadows2008}
Meadows, Donella~H.
\newblock \emph{Thinking in Systems: A Primer}.
\newblock Edited by Diana Wright.
\newblock Chelsea Green Publishing, White River Junction, VT, 2008.

\end{thebibliography}

\end{document}
